\chapter{Beyond Prompting: The Shift from Extraction to Collaboration}

The dominant paradigm for interacting with LLMs treats them as sophisticated search engines. You ask a question; you get an answer. The better your question, the better your answer. This is \textbf{extraction-based thinking}.

But extraction misses the deeper capability of these systems. LLMs don't merely retrieve---they \emph{generate}. They don't just recall patterns---they \emph{synthesize}. And synthesis, when properly guided, can produce something neither you nor the model possessed alone.

\section{The Fundamental Shift}

\textbf{DCF Principle \#1: Collaboration Over Extraction}

\begin{table}[H]
\centering
\begin{tabular}{ll}
\toprule
\textbf{Extraction Paradigm} & \textbf{DCF Collaboration Paradigm} \\
\midrule
``Give me the answer'' & ``Help me think through this'' \\
One-shot prompts & Recursive dialogue \\
Evaluation: correctness & Evaluation: clarity gained \\
User as interrogator & User as co-creator \\
Output is the goal & Understanding is the goal \\
\bottomrule
\end{tabular}
\caption{Extraction vs. Collaboration Paradigms}
\end{table}

This shift transforms the LLM from a tool into a \textbf{thinking partner}---a system that scaffolds cognition rather than replacing it.

\begin{quotebox}
\textbf{Agentic Era Update:} In systems like Claude Code, this shift is partially automated. The agent \emph{already} engages in multi-step reasoning, plans before acting, and self-corrects. Your role shifts from orchestrating collaboration to \emph{evaluating} whether the agent's autonomous collaboration with itself produced the right outcome. DCF still applies---at the review and approval stages.
\end{quotebox}
